\section{Background and Related Work}
\label{sec:related}
%% Describe the following papers
%
%[1] J. W. Vaupel and A. I. Yashin. Heterogeneity's ruses: Some surprising effects of selection on population dynamics. The American Statistician 39(3), pp. 176-185. 1985.
%
%[2] P. J. Bickel, E. A. Hammel and O'connell J.W. Sex bias in graduate admissions: Data from berkeley. Science (New York, N.Y.) 187(4175), pp. 398-404. 1975.
%
%[3] P. Singer et al. Evidence of online performance deterioration in user sessions on reddit. PLoS One 11(8), 2016.
%
%[4] S. Barbosa et al. Averaging gone wrong: Using time-aware analyses to better understand behavior. 2016. . DOI: 10.1145/2872427.2883083.
%
%[5] E. Ferrara et al. Dynamics of Content Quality in Collaborative Knowledge Production. ICWSM 2017.

%Ecological fallacy
%
%1957 paper???
%
%% There was another paper Nazanin found showing how to detect Simpson's paradox.
%
%Miguel A. Hernan, David Clayton and Niels Keiding (2011) Simpson's paradox unraveled. Int. J Epidemiology 40(3): 780-785
%"Key messages: Identical data arising from different causal structures (i.e., Simpson's paradox) needs to be analyzed differently. Causal analyses need to be guided by subject matter knowledge. No purely statistical rules exist to guide causal analyses."

The goal of data analysis is to identify important associations between features, or variables, in data. %covariates
However, hidden correlations between variables can lead analysis to wrong conclusions. One important manifestation of this effect is Simpson's paradox, according to which \emph{an association that appears in different subgroups of data may disappear, and even reverse itself, when data is aggregated across subgroups}.
Instances of the paradox have been documented across a variety of disciplines, including demographics, economics, political science, and clinical research, and it has been argued that the presence of Simpson's paradox implies that interesting patterns exist in data~\cite{fabris2000discovering}.
A notorious example of Simpson's paradox  arose during a gender bias lawsuit against UC Berkeley. Analysis of graduate school admissions data revealed a statistically significant bias against women. %However, when admissions data was disaggregated by department, the bias disappeared, and in some departments, women even had a slight advantage over men.
However, the pattern of discrimination observed in this aggregate data disappeared when admissions data was disaggregated by department.
Bickel et al.~\cite{Bickel1975} attributed this effect to Simpson's paradox. They argued that the subtle correlations between the popularity of departments among the genders and their selectivity resulted in women applying to departments that were hardest to get into, which skewed analysis.

Simpson's paradox must also be considered in the analysis of trends. Vaupel and Yashin~\cite{Vaupel85heterogeneity} give a compelling illustration of how survivor bias can shift the composition of data, distorting the conclusions drawn from it. Analysis of recidivism among convicts released from prison shows that the rate at which they return to prison declines over time. From this, policy makers may conclude that age has a pacifying effect on crime: older convicts are less likely to commit crimes. In reality, this conclusion is false. %and there is no evidence that older convicts are less likely to re-offend.
Instead, we can think of the population of ex-convicts as composed of two subgroups with constant, but very different recidivism rates. The first subgroup, let's call them ``reformed,'' will never commit a crime once released from prison. The other subgroup, the ``incorrigibles,'' will always commit a crime. Over time, as ``incorrigibles'' commit offenses and return to prison, there are fewer of them left in the population. The survivor bias changes the composition of the population under study, creating an illusion of an overall decline in recidivism rates. As Vaupel and Yashin warn, ``unsuspecting researchers who are not wary of heterogeneity's ruses may fallaciously assume that observed patterns for the population as a whole also hold on the sub-population or individual level.'' Their paper gives numerous other examples of such \emph{ecological fallacies}.

Similar illusions crop up in many studies of social behavior. For example, when examining how social media users respond to information from their friends (other users that they follow), it may appear that if \emph{more} of a user's friends use a hashtag then the user will be \emph{less} likely to use it himself or herself~\cite{Romero11www}. Similarly, the more friends share some information, the less likely the user is to share it with his or her followers~\cite{Versteeg11icwsm}. From this, one may conclude the additional exposures to information in a sense ``innoculate'' the user and act to suppress the sharing of information. In fact, this is not the case, and instead, additional exposures monotonically increase the user's likelihood to share information with followers~\cite{Lerman2016futureinternet}. However, those users who follow many others, and are likely to be exposed to information or a hashtag multiple times, are less responsive overall, because they are overloaded with information they receive from all the friends they follow~\cite{Hodas12socialcom}. Calculating response as a function of the number of exposures in the aggregate data falls prey to survivor bias: the more responsive users (with fewer friends) quickly drop out of the average (since they are generally exposed fewer times), leaving the highly connected, but less responsive, users behind.  The reduced susceptibility of these highly connected users biases aggregate response, leading to wrong conclusions about individual behavior. Once data is disaggregated based on the volume of information individuals receive, a clearer pattern of response emerges, one that is more predictive of behavior~\cite{Hodas14srep}. Multiple examples of Simpson's paradox have been identified in empirical studies of online behavior. A study~\cite{Barbosa2016} of Reddit found that while it may appear that average comment length on decreases over any fixed period of time, when data is disaggregated into groups based on the year user joined Reddit, comment length within each group increases during the same time period. It is only because users who joined early tend to write longer comments that the Simpson's paradox appears.

Data heterogeneity also impacts statistical analysis of data~\cite{estes1956problem} and causal inference~\cite{Xie2013population}. However, no general framework to recognize and mitigate Simpson's paradox exist. Current methods require that the structure of data be explicitly specified~\cite{Bareinboim2016causal} or at best be guided by subject matter knowledge~\cite{Hernan2011}.

%Miguel A. Hernan, David Clayton and Niels Keiding (2011) Simpson's paradox unraveled. Int. J Epidemiology 40(3): 780-785
%"Key messages: Identical data arising from different causal structures (i.e., Simpson's paradox) needs to be analyzed differently. Causal analyses need to be guided by subject matter knowledge. No purely statistical rules exist to guide causal analyses."

Despite accumulating evidence that Simpson's paradox affects inference from data~\cite{Xie2013population,estes1956problem}, scientists do not routinely test for the presence of this paradox in heterogeneous data. %Our work addresses this problem by proposing a framework for testing when trends found in aggregate data could arise due to Simpson's paradox.
Our work addresses this knowledge gap by proposing a statistical method to systematically uncover instances of Simpson's paradox in data.






